%; whizzy paragraph -pdf xpdf -latex ./whizzypdfptex.sh
%; whizzy-paragraph "^\\\\begin{frame}\\|\\\\emtext"
% latex beamer presentation.
% platex, latex-beamer でコンパイルすることを想定。 

%     Tokyo Debian Meeting resources
%     Copyright (C) 2012 Junichi Uekawa

%     This program is free software; you can redistribute it and/or modify
%     it under the terms of the GNU General Public License as published by
%     the Free Software Foundation; either version 2 of the License, or
%     (at your option) any later version.

%     This program is distributed in the hope that it will be useful,
%     but WITHOUT ANY WARRANTY; without even the implied warreanty of
%     MERCHANTABILITY or FITNESS FOR A PARTICULAR PURPOSE.  See the
%     GNU General Public License for more details.

%     You should have received a copy of the GNU General Public License
%     along with this program; if not, write to the Free Software
%     Foundation, Inc., 51 Franklin St, Fifth Floor, Boston, MA  02110-1301 USA

\documentclass[cjk,dvipdfmx,12pt]{beamer}
\usetheme{Tokyo}
\usepackage{monthlypresentation}

%  preview (shell-command (concat "evince " (replace-regexp-in-string "tex$" "pdf"(buffer-file-name)) "&")) 
%  presentation (shell-command (concat "xpdf -fullscreen " (replace-regexp-in-string "tex$" "pdf"(buffer-file-name)) "&"))
%  presentation (shell-command (concat "evince " (replace-regexp-in-string "tex$" "pdf"(buffer-file-name)) "&"))

%http://www.naney.org/diki/dk/hyperref.html
%日本語EUC系環境の時
\AtBeginDvi{\special{pdf:tounicode EUC-UCS2}}
%シフトJIS系環境の時
%\AtBeginDvi{\special{pdf:tounicode 90ms-RKSJ-UCS2}}

\newenvironment{commandlinesmall}%
{\VerbatimEnvironment
  \begin{Sbox}\begin{minipage}{1.0\hsize}\begin{fontsize}{8}{8} \begin{BVerbatim}}%
{\end{BVerbatim}\end{fontsize}\end{minipage}\end{Sbox}
  \setlength{\fboxsep}{8pt}
% start on a new paragraph

\vspace{6pt}% skip before
\fcolorbox{dancerdarkblue}{dancerlightblue}{\TheSbox}

\vspace{6pt}% skip after
}
%end of commandlinesmall

\title{東京エリアDebian勉強会}
\subtitle{第112回 2014年4月度}
\author{野島貴英}
\date{2014年4月19日}
\logo{\includegraphics[width=8cm]{image200607/openlogo-light.eps}}

\begin{document}

\begin{frame}
\titlepage{}
\end{frame}

\begin{frame}{設営準備にご協力ください。}
会場設営よろしくおねがいします。
\end{frame}

\begin{frame}{Agenda}
 \begin{minipage}[t]{0.45\hsize}
  \begin{itemize}
   \item 注意事項
	 \begin{itemize}
	  \item 写真はセミナールーム内のみ可です。
          \item 出入りは自由でないので、もし外出したい方は、野島まで一声くださいませ。
	 \end{itemize}
   \item 最近あったDebian関連のイベント報告
	 \begin{itemize}
	  \item 第111回 東京エリアDebian勉強会
	 \end{itemize}
  \end{itemize}
 \end{minipage} 
 \begin{minipage}[t]{0.45\hsize}
  \begin{itemize}
   \item Debian Trivia Quiz
  \end{itemize}
 \end{minipage}
\end{frame}

\section{イベント報告}
\emtext{イベント報告}

\begin{frame}{第111回 東京エリアDebian勉強会}

 東京エリアDebian勉強会111回目は(株)スクウェア・エニックスさんで開催されました。
5名の参加者がありました。

\begin{itemize}
\item debianからiphone5(iOS7.1)につなぐ件について、
  \begin{itemize}
    \item debianパッケージを作ってつなぎ、
    \item 公開情報に基づいた接続技術
　\end{itemize}
について発表がありました。
\item 参加者全員で、各自の作業を行い、最後に成果発表をしました。
\end{itemize}

\end{frame}

\section{Debian Trivia Quiz}
\emtext{Debian Trivia Quiz}
\begin{frame}{Debian Trivia Quiz}

  Debian の常識、もちろん知ってますよね?
知らないなんて恥ずかしくて、知らないとは言えないあんなことやこんなこと、
みんなで確認してみましょう。

今回の出題範囲は\url{debian-devel-announce@lists.debian.org},
\url{debian-devel@lists.debian.org} に投稿された
内容などからです。

\end{frame}

\subsection{問題}

%; whizzy-master ../debianmeetingresume201311.tex
% $B0J>e$N@_Dj$r$7$F$$$k$?$a!"$3$N%U%!%$%k$G(B M-x whizzytex $B$9$k$H!"(Bwhizzytex$B$,MxMQ$G$-$^$9!#(B
%

\santaku
{2015$BG/$N(BDebconf15$B$N3+:E9q$O$I$3$K$J$C$?$G$7$g$&!)(B}
{$B%\%9%K%"!&%X%k%D%'%4%S%J(B}
{$B%&%/%i%$%J(B}
{$B%I%$%D(B}
{C}
{$B:#EY$O%I%$%D$@$=$&$G$9!#$A$J$_$K:#G/$N(BDebconf14$B$O(B8/23-31$B$G(BUSA$B$N(BPortland,Oregon
$B$G3+:EM=Dj$G$9!#(B}

\santaku
{DPL$BA*5s$,9T$o$l$^$7$?!#:#G/$N(BDPL$B$OC/$K$J$C$?$G$7$g$&$+!)(B}
{Lucas Nussbaum}
{Neil McGovern}
{Rapha\"{e}l Hertzog}
{A}
{$B:#G/$N(BDPL$B$N8uJd<T$O#2L>$G!"(BLucas Nussbaum$B$5$s!"(BNeil McGovern$B$5$s$N#2L>(B
$B$G$7$?!##2L>$H$b<+A&$H$N$3$H$G$9!#A*5s$N7k2L!"(B
Lucas Nussbaum(lucus)$B$5$s$N05>!$@$C$?$h$&$G$9!#$A$J$_$K!"(B
Rapha\"{e}l Hertzog$B$5$s$O!"(BThe Debian 
Administrator\'s handbook$B$N:n<T!"B>$K$b0N6H$,$?$/$5$s!#(B}

\santaku
{clang3.4$B$K$h$k(BDebian$B%Q%C%1!<%8$N:F9=C[$,9T$o$l$^$7$?!#7k2L2?(B\%$B$N%Q%C%1!<%8$,@.8y$7$?$G$7$g$&$+!)(B}
{90\%}
{50\%}
{10\%}
{A}
{\url{http://clang.debian.net/}$B$,(Bclang$B$K$h$k(BDebian$B%Q%C%1!<%8:F9=C[$N%]!<%?%k%5%$%H$G$9!#KhG/!"(Bclang$B$N%P!<%8%g%s$r>e$2$F!"A4(BDebian$B%Q%C%1!<%8$r:F9=C[$7$?7k2L$,:\$j$^$9!#:#G/$N7k2L$H$7$F$O:F9=C[BP>]$N%Q%C%1!<%8$N?t$,5nG/$HBgI}$KA}$($F$$$k$K$b$+$+$o$i$:!"9=C[<:GT$K=*$o$C$?%Q%1!<%8?t$,5nG/$HJQ$o$i$J$+$C$?$H$$$&Hs>o$KNI$$7k2L$H$J$C$F$$$^$9!#(B}

\santaku
{beagle board$B%7%j!<%:$H$$$&Hs>o$K?M5$$N$"$k(BARM$B$N<B83%\!<%I$K%P%s%I%k$5$l$k(BOS$B$N>-Mh$N8+DL$7$K$D$$$F!"(Bbeagle board$B$NAON)<T$,$I$N(BOS$B$K$9$kM=Dj$HH/8@$7$?$+!)(B}
{Gentoo}
{Debian$B$C$7$g!*(BDebian}
{Andoroid OS}
{B}
{\url{http://opensource.com/life/14/3/interview-jason-kridner-beagleboard}
$B$K$F!">-Mh(BDebian$B$K$9$k$H$$$&H/8@$,$"$j$^$9!#$H$3$m$G!"(Bbeagle board$B%7%j!<%:$OL$$@$K(BOMAP$B$r(B
$B;H$$B3$1$k$N$+$,6=L#DE!9$G$O$"$j$^$9!#(B}

\santaku
{$B7cO@$NKv!"(B3$B7nCf=\:"$K(BDebian$B$N(Bca-certificates$B%Q%C%1!<%8$+$i>C$($?(Broot$B>ZL@=q$,$"$j$^$9!#$=$l$O$J$s$G$7$g$&!)(B}
{RapidSSL$B$N(Broot$B>ZL@=q(B}
{CAcert$B$N(Broot$B>ZL@=q(B}
{Verisign$B$N(Broot$B>ZL@=q(B}
{B}
{$B5DO@$N%5%^%j$O!"(BLWN$B$N5-;v(B\url{https://lwn.net/Articles/590879/}$B$,H=$j$d$9$$$G$9!#$^$?!"(BCAcert$B$C$F2?!)$H$$$&J}$O!"Bh(B71$B2sEl5~%(%j%"(BDebian$BJY6/2q(B(2010$BG/(B12$B7n3+:E(B)\url{http://tokyodebian.alioth.debian.org/2010-12.html}$B$K7G:\$5$l$F$$$kJY6/2q;qNA$,$*4+$a$G$9!#(B}

\santaku
{Jessie$B$K$F%G%9%/%H%C%W4D6-$rA*Br$7$?:]$KF3F~$5$l$k!"%G%U%)%k%H%3%_%e%K%1!<%7%g%s%D!<%k$N8uJd$K$D$$$F5DO@$,$5$l$F$$$j$^$9!#0J2<$N$I$l$G$7$g$&!)(B}
{Empathy}
{licq}
{jitsi}
{C}
{Debian$B$G$O!"%G%U%)%k%H$N%3%_%e%K%1!<%7%g%s%D!<%k$H$7$F!"$[$\40A4$K(BRTC/VoIP$B$r%5%]!<%H$9$k$3$H$,K>$^$7$$$H$5$l$?$?$a!"$3$A$i$K8~$$$F$$$k%D!<%k$H$7$F:#$^$G$N(BEmpathy$B$h$j$b(Bjitsi$B$,8~$$$F$$$k$N$G$O!)$H$$$&$3$H$+$i5DO@$,3+;O$5$l$^$7$?!#(B}

\santaku
{$B@hF|(BDebian$B$K(BOTR$B%A!<%`$H$$$&(BOTR$B%=%U%H$r%Q%C%1!<%82=$9$k%0%k!<%W$,7k@.$5$l$^$7$?!#$H$3$m$G(BOTR$B$C$F$J$s$NN,!)(B}
{Owa-Tte-Ru}
{OpticalTRacking}
{Off-the-Record}
{C}
{Off-the-Record$B$H$O!"0E9f2=5;=Q$r;H$C$F%$%s%U%iDs6!<T$K$9$i%a%C%;!<%8$N$d$j$H$j$NFbMF$r8+$;$J$$!J5-O?$5$;$J$$!K%a%C%;!<%8%5!<%S%9$rL\;X$7$?$b$N$G$9(B\url{https://www.otr.im/}$B!#(B}

\santaku
{$B@hF|!"(BDebian squeeze$B$N%5%]!<%H4|4V$,?-$S$k@k8@$,(BDSA$B$K$h$j%"%J%&%s%9$5$l$^$7$?!#7k6I$$$D$K$J$C$?!)(B}
{2015/2}
{2016/2}
{2016/5}
{B}
{Long Term Support(LTS)$B$@$=$&$G$9(B\url{https://lists.debian.org/debian-security-announce/2014/msg00082.html}$B!#M=Dj$G$O!"(B2014/5/31$B:"$K%5%]!<%H=*N;$9$k$O$:$G$7$?$N$G!"#2G/<e$N1dD9$H$J$j$^$9!#$?$@%5%]!<%H1dD9$5$l$k$N$O!"(BDebian squeeze$B$NA4It$N%Q%C%1!<%8$G$O$J$$$N$G!"%5%]!<%H$5$l$J$$%Q%C%1!<%8$r==J,$K$*;H$$$N3'MM$O!"(BDebian wheezy$B$X%"%C%W%0%l!<%I$9$k$3$H$r$*4+$a$7$F$*$-$^$9!#(B}




\section{事前課題}
\emtext{事前課題}
{\footnotesize
\begin{prework}{ 野島貴英 }

\begin{itemize}
\item Xmrisパッケージ化頑張る
\item iOS7.1 debian接続再び。
\end{itemize}
\end{prework}

\begin{prework}{ 岩松 信洋 }

?

\end{prework}

\begin{prework}{ dictoss(杉本　典充) }

Debian上でLeap Motionを動かしてみる

\end{prework}

\begin{prework}{ Yoshida Shin }

GitHub に公開しているプログラムの機能追加。
動作環境は Debian Sid で。
\end{prework}

\begin{prework}{ Koji Hasebe }

Debianパッケージ作成の方法を調査し、実際に作成する。
アウトプットとして、その方法を手順書にまとめる。
\end{prework}

\begin{prework}{ henrich }
\begin{itemize}
\item OSC北海道用の資料作成
\item debhelperの査読の続き
\item パッケージまわりでのなにか
\end{itemize}
のいずれか。

\end{prework}

\begin{prework}{ zinrai }

Debianパッケージ作成。
\begin{itemize}
\item Debianパッケージ作成してみたい
\item 「Debianパッケージング道場」に都合がつかず参加できなかったので「もくもく会」の時間を使いパッケージ作成の作法を心得ている人がいる環境でパッケージ作成したい
\item redis 2.8.8をDebianパッケージ作成の題材にしてみようと思う。
\url{http://redis.io/}
\item 通勤時間などを使い当日までにこの辺に目を通しておきたい
\url{http://www.debian.or.jp/community/devel/}
\end{itemize}
\end{prework}

\begin{prework}{ yyuu }

自分で管理してるプロジェクトのdeb化を要望されてしまったので、そのための作業を行ないたいと考えています。

\url{https://github.com/yyuu/pyenv/issues/138}

\end{prework}

\begin{prework}{ 野首(@knok) }
\begin{itemize}
\item KAKASIローマ字テーブルの改善
\item LanguageToolのルール追加
\item gnu.org, veの翻訳
\end{itemize}
\end{prework}

\begin{prework}{ まえだこうへい }

\begin{itemize}
\item Golang関係のパッケージ化の話をする
\item \url{http://qa.debian.org/developer.php?login=mkouhei@palmtb.net}
のバグ潰し＆パッケージアップデート
\item 上記のGolang関係のパッケージをpkg-go teamでのメンテナンスを始める（MLに投稿＆ITPする）
\end{itemize}
\end{prework}

}

\section{Serf}
\emtext{Serf}

\section{今後のイベント}
\emtext{今後のイベント}
\begin{frame}{今後のイベント}
\begin{itemize}
 \item 2014年5月17日(土) 東京エリアDebian勉強会
\end{itemize}
\end{frame}

\section{今日の宴会場所}
\emtext{今日の宴会場所}
\begin{frame}{今日の宴会場所}
未定
\end{frame}

\end{document}

;;; Local Variables: ***
;;; outline-regexp: "\\([ 	]*\\\\\\(documentstyle\\|documentclass\\|emtext\\|section\\|begin{frame}\\)\\*?[ 	]*[[{]\\|[]+\\)" ***
;;; End: ***
